% SOURCE_AUTHORITY: sga5_fr_workpass.tex, lines 15076--15102 (LF-normalized unit).
% SOURCE_SHA256: 791F4EFFC5E02832D5D77ED03518C8156D6F07E4C8238B03545DB93D883FBB28
\subsubsection*{n\textsuperscript{o} 3. Demonstração do teorema}

No final dessas reduções, só resta estabelecer o teorema quando $X$ é a reta projetiva. Primeiro, transformaremos a fórmula (1) do enunciado em uma fórmula de tipo Lefschetz relativa a traços; esta última fórmula será demonstrada no caso particular em que $X$ seja uma curva projetiva e lisa.

Vamos retomar a fórmula (1):
\[
L_F=\prod_i\bigl(L_{\R^i_!g(F)}\bigr)^{(-1)^i}
=\frac{P_F^1\cdots P_F^{2n-1}}{P_F^0\cdots P_F^{2n}},
\quad\text{com}\quad
P_F^i(t)=\det(1-f^{-1}_{H^i_!(\bar X,\bar F)}t).
\]
O logaritmo do primeiro membro escreve-se:
\[
\sum_{x\in X^0}\sum_{n=1}^{\infty}\Tr(f_{F_{\bar x}}^{-n d(x)})\,t^{n d(x)}/n
=\sum_{n=1}^{\infty}\left(\sum_{d(x)|n} d(x)\Tr(f_{F_{\bar x}}^{-n})\right)t^n/n,
\]
enquanto o do segundo membro é
\[
\sum_{n\ge1}\left(\sum_i(-1)^i\Tr(f^{-n}_{H^i_!(\bar X,\bar F)})\right)t^n/n.
\]
Comparando os coeficientes de $t^n/n$, vê-se que é preciso estabelecer:
\[
\sum_{d(x)|n}d(x)\Tr(f_{F_{\bar x}}^{-n})
=\sum_i(-1)^i\Tr(f^{-n}_{H^i_!(\bar X,\bar F)})
\tag{2}
\]
para todo $n\ge1$; estes cálculos são legítimos porque o corpo $\Omega$ tem característica zero.
